\clearpage\SourceLeaf{436}\NoSourceHead
\HistoricTitle{BRIEFWECHSEL.}
\NumberedItem{1}
\HistoricTitle{\textsc{Gauss} an C. F. \textsc{Hindenburg}.}
\begin{flushright}Braunschweig d. 8\textsuperscript{ten} Oct. 1799.\end{flushright}
\noindent\hspace*{.18\linewidth}Wohlgeborener\par
\noindent\hspace*{.18\linewidth}Hoch zu verehrender Herr Professor.

\indent Ich nehme mir die Freiheit, Ew. Wohlgeboren hiemit ein Exemplar einer vor Kurzem von mir fertig gewordenen kleinen Schrift[\textsuperscript{*})] zukommen zu lassen, und bitte dasselbe als ein geringes Zeichen meiner Verehrung von Dero grossen Verdiensten gütigst anzunehmen. Vielleicht haben Dieselben auf eine oder die andere Art Gelegenheit zum allgemeinen Bekanntwerden dieser Schrift etwas beizutragen; mir aber wird es die grösste Freude sein, wenn Ew. Wohlgeboren derselben den schätzbarsten Beifall gönnen wollen. Erlauben Dieselben mir nur die eine Bemerkung, dass die Kürze, Gedrungenheit, womit sie geschrieben ist, in mehr als Einer sehr wichtigen Ursache ihren Grund hat; zwar hat das die Folge, dass nicht jeder gleich beim ersten flüchtigen Durchlesen sich des Ganzen bemächtigt; allein diess hat bei einem Aufsatze, der eigentlich nur für die Geübtere bestimmt sein kann, Nichts auf sich, und ich habe schon mehrere überzeugende Belege, dass man bei aufmerksamer nicht übereilter Lectüre keine Dunkelheit darin gefunden hat.

\indent Es sind schon drittehalb Jahre, dass ich mir einmal die Freiheit nahm Ew. Wohlgeboren zu schreiben und Denenselben einen kleinen Aufsatz über
\par\medskip\noindent\rule{.24\linewidth}{.4pt}\par\smallskip
{\small
\noindent\hspace*{\parindent}[\textsuperscript{*}) \textit{Demonstratio nova etc.} Helmstadii 1799, Werke III, S. 1.]
\par}
\clearpage\SourceLeaf{437}\ArithmeticHead{430}{L}{ANALYSIS. BRIEFWECHSEL.}
\noindent das wichtige Theorem \textsc{von Lagrange} zu communiciren[\textsuperscript{*})], wovon Hr. H[of] R[ath] \textsc{Kästner} in Göttingen damals die Besorgung gütigst übernehmen wollen. Da ich nicht gefunden, dass Ew. Wohlgeboren davon in Dero \textls[120]{Archiv} Gebrauch gemacht und ich auch von Denenselben auf mein Schreiben keine Antwort erhalten habe, so muss ich fast daraus schliessen, dass dasselbe Ihnen nicht zu Händen gekommen sei. Dem sei indess wie ihm wolle, so ist es mir gegenwärtig aus mehreren Ursachen lieb, dass mein Beweis jenes Lehrsatzes \textls[120]{nicht} gedruckt ist, theils weil ich das \textsc{Laplace}sche Verfahren[\textsuperscript{**})], welches mir damals noch unbekannt war, weit vorziehe und für eleganter und subtiler halte, theils weil ich nach der Hand zu dem Theorem selbst erhebliche Erweiterungen gefunden habe[\textsuperscript{***})], welche ich bekannt machen werde, sobald andere Arbeiten mir Musse dazu lassen. Indessen würde es mir lieb sein wenn Ew. Wohlgeboren mich zu unterrichten die Güte haben wollten, ob Dieselben gedachten Aufsatz erhalten haben oder nicht, und im ersten Fall mir denselben wieder zurückschicken, zumal da Dieselben nach der ausführlichen Bearbeitung dieser Gegenstände von meinem verehrungswürdigsten Freunde, Herrn Prof. \textsc{Pfaff}[\textsuperscript{\ensuremath{\dagger}})] ohnehin keinen Gedanken mehr haben könnten, Gebrauch davon zu machen.

\indent Von meinen in voriger Messe angekündigten \textit{Disquiss. Arithmett.} ist, leider, der Druck durch verschiedene verdriessliche Umstände sehr verzögert; das Werk wird etwa zwischen 30 und 40 Bogen stark werden, wovon kaum die Hälfte abgedruckt ist; indessen hoffe ich den Druck, welcher seit einigen Monaten ganz unterbrochen gewesen ist, bald wieder in raschen Gang zu bringen und das Ganze dann bald zu vollenden.
\par\medskip\noindent\rule{.24\linewidth}{.4pt}\par\smallskip
{\small
\noindent\hspace*{\parindent}[\textsuperscript{*}) Dieser Aufsatz ist nach einer im Nachlaß (Fb, Kapsel 46a) befindlichen Handschrift abgedruckt Werke VIII, S. 76; vergl. auch die Tagebuchaufzeichnung Nr. 49 vom 27. Dezember 1796.]
\par\noindent\hspace*{\parindent}[\textsuperscript{**}) Siehe den VII. Abschnitt der Abhandlung von \textsc{Laplace} \textit{Mémoire sur l’usage du calcul aux différences partielles dans la théorie des suites}, Histoire de l’Acad. des Sciences, Année 1777, Paris 1780, Mémoires etc. S. 113; vergl. \textit{Mécanique céleste} I, Paris, an VII (1798—99), livre II, Nr. 22.]
\par\noindent\hspace*{\parindent}[\textsuperscript{***}) Siehe die Tagebuchaufzeichnung Nr. 86 vom Mai 1798.]
\par\noindent\hspace*{\parindent}[\textsuperscript{\ensuremath{\dagger}}) Außer der Arbeit von \textsc{J. Fr. Pfaff}, die \textsc{Gauss} in dem \textsc{Hindenburg} eingereichten Aufsatze anführt, (siehe Werke VIII, S. 76 letzte Zeile) \textit{Analysis einer wichtigen Aufgabe des Herrn de la Grange}, Archiv der r. u. a. Mathematik \textbf{1}, Heft 1, S. 81, kommen noch in Betracht die Aufsätze \textit{Allgemeine Summation einer Reihe, worin höhere Differenziale vorkommen}, ebenda \textbf{1}, Heft 3, 1795, S. 337 und \textbf{2}, Heft 5, 1796, S. 67, ferner \textit{Auszug aus einem Briefe vom Herrn Prof. Pfaff an den Herausgeber}, ebenda \textbf{2}, Heft 7, 1797, S. 347 und \textit{Disquisitiones analyticae}, Helmstadii 1797, 3. Abhandlung: \textit{Investigatio serierum transcendenter summabilium}, S. 65—132.]
\par}
\clearpage\SourceLeaf{438}\ArithmeticHead{431}{R}{GAUSS AN SCHUMACHER.}
\indent Mit der vollkommensten Verehrung habe ich die Ehre zu beharren Ew. Wohlgeboren
\begin{flushright}ganz ergebenster Diener\par Doctor \textsc{Gauss}.\end{flushright}
\NumberedItem{2}
\HistoricTitle{\textsc{Gauss} an \textsc{Schumacher}.}
\HistoricCitation{[Briefwechsel zwischen \textsc{Gauss} und \textsc{Schumacher} VI, Altona 1865, S. 51.]}
\noindent {\ensuremath{{.}\;{.}\;{.}\;{.}}} Über {\ensuremath{{.}\;{.}\;{.}}} die Convergenz der Reihe für die Mittelpunktsgleichung, kann ich für jetzt mich nicht weiter auslassen. Da ich, lange vor 1817 (wo ich \textsc{Carlinis} Abh[andlung] in den Effem. für 1818[\textsuperscript{*})] erhielt) die Aufgabe selbst auf eine ohne allen Vergleich kürzere Art aufgelöset hatte, so habe ich damals diese Abhandl[ung] ebenso wie jetzt \textsc{Jacobis} Aufsatz[\textsuperscript{**})] nur ganz flüchtig angesehen, und nachdem ich in jener die Übereinstimmung des Hauptresultats mit dem meinigen bemerkt hatte, nicht weiter gelesen; daher war das von \textsc{Jacobi} jetzt gerügte Versehen von mir nicht bemerkt.
\begin{flushright}Stets der Ihrige\hspace*{.18\linewidth}\par C. F. \textsc{Gauss}.\end{flushright}
\indent Göttingen, 4. Dec[ember] 1849.
\par\medskip\noindent\rule{.24\linewidth}{.4pt}\par\smallskip
{\small
\noindent\hspace*{\parindent}[\textsuperscript{*}) \textit{Ricerche sulla convergenza della serie che serve alla soluzione del problema di Keplero}, Memoria di \textsc{Francesco Carlini}, Milano 1817; auch als \textls[120]{Appendice} zu den Effemeridi di Milano 1818.]
\par\noindent\hspace*{\parindent}[\textsuperscript{**}) \textit{Über die annähernde Bestimmung sehr entfernter Glieder in der Entwickelung der elliptischen Coordinaten, nebst einer Ausdehnung der Laplaceschen Methode zur Bestimmung der Functionen grosser Zahlen}, Astronomische Nachrichten \textbf{28}, Nr. 665, März 1849, S. 257, \textsc{Jacobis} Werke VII, S. 175.]
\par}
